(1) 对 \(b/v\ll1\) 展开 \(1/(v-b)=v^{-1}+bv^{-2}+\cdots\)，得到 \(p=RT/v+(RTb-a)/v^2+\cdots\)。迭代反演为 \(v=RT/p+b-a/(RT)+O(p)\)。括号内第二维里系数 \(B_2(T)=b-a/(RT)\)，分别对应短程排斥体积及平均吸引能。

(2) 焓的 Maxwell 微分恒等式为 \(dH=C_p\,dT+[v-T(\partial v/\partial T)_p]dp\)。对节流流体 \(dH=0\)，因而 \(\mu_{\rm JT}=(\partial T/\partial p)_H=[T(\partial v/\partial T)_p-v]/C_p\)。节流不是等熵可逆过程，求系数仅用入口出口焓相等，而不可另加错误的恒熵条件。

(3) 首阶近似的导数是 \((\partial v/\partial T)_p=R/p+a/(RT^2)+O(p)\)，代回得 \(\mu_{\rm JT}\simeq[2a/(RT)-b]/C_p\)。反转温度 \(T_{\rm inv}\simeq2a/(Rb)\)。低于此温度 \(\mu_{\rm JT}>0\)，节流时 \(dp<0\) 所以 \(dT<0\) 冷却；高于此温度则节流升温。

(4) 在理想极限 \(a=b=0\)，有 \(v=RT/p\)，恒焓过程保持温度、\(\mu_{\rm JT}=0\)。分子 \(a/(RT)\) 与 \(b\) 都有摩尔体积量纲，除以 \(C_p\) 得温度每压强，符合定义。反转温度只是低压极限值；接近凝聚、临界或 \(p\) 不足够小的区域应回到完整状态方程。
